\documentclass[10pt,a4paper]{amsart}
\usepackage{amsmath,amssymb,amsthm,graphicx}
% Retained preamble for the Stacks Project chapter snapshot in this directory.
% The chapter's own preamble.tex is a copy of the shared Stacks Project preamble; the block
% below is sliced MECHANICALLY, line for line, from it. Nothing was added, removed or
% redefined; the wrapper class, the xr-hyper/hyperref wiring and the externaldocument list
% are omitted because the assembler supplies the document class.
%
% lines taken from preamble.tex: 3, 5, 9, 11, 12, 156, 157, 158, ... (51 lines total)
%
% package coverage restored from the chapter's own preamble (2026-10-13): xy, xy-2cell

\usepackage[all]{xy}\usepackage{verbatim}
\xyoption{2cell}
\UseAllTwocells

\newenvironment{reference}{\comment}{\endcomment}
\newenvironment{slogan}{\comment}{\endcomment}
\newenvironment{history}{\comment}{\endcomment}
\theoremstyle{plain}
\newtheorem{theorem}[subsection]{Theorem}
\newtheorem{proposition}[subsection]{Proposition}
\newtheorem{lemma}[subsection]{Lemma}
\theoremstyle{definition}
\newtheorem{definition}[subsection]{Definition}
\newtheorem{example}[subsection]{Example}
\newtheorem{exercise}[subsection]{Exercise}
\newtheorem{situation}[subsection]{Situation}
\theoremstyle{remark}
\newtheorem{remark}[subsection]{Remark}
\newtheorem{remarks}[subsection]{Remarks}
\numberwithin{equation}{subsection}
\def\lim{\mathop{\mathrm{lim}}\nolimits}
\def\colim{\mathop{\mathrm{colim}}\nolimits}
\def\Spec{\mathop{\mathrm{Spec}}}
\def\Hom{\mathop{\mathrm{Hom}}\nolimits}
\def\Ext{\mathop{\mathrm{Ext}}\nolimits}
\def\SheafHom{\mathop{\mathcal{H}\!\mathit{om}}\nolimits}
\def\SheafExt{\mathop{\mathcal{E}\!\mathit{xt}}\nolimits}
\def\Sch{\mathit{Sch}}
\def\Mor{\mathop{\mathrm{Mor}}\nolimits}
\def\Ob{\mathop{\mathrm{Ob}}\nolimits}
\def\Sh{\mathop{\mathit{Sh}}\nolimits}
\def\NL{\mathop{N\!L}\nolimits}
\def\CH{\mathop{\mathrm{CH}}\nolimits}
\def\proetale{{pro\text{-}\acute{e}tale}}
\def\etale{{\acute{e}tale}}
\def\QCoh{\mathit{QCoh}}
\def\Ker{\mathop{\mathrm{Ker}}}
\def\Im{\mathop{\mathrm{Im}}}
\def\Coker{\mathop{\mathrm{Coker}}}
\def\Coim{\mathop{\mathrm{Coim}}}
\DeclareMathSymbol{\boxtimes}{\mathbin}{AMSa}{"02}
\def\QCohstack{\mathcal{QC}\!\mathit{oh}}
\def\Cohstack{\mathcal{C}\!\mathit{oh}}
\def\Spacesstack{\mathcal{S}\!\mathit{paces}}
\def\Quotfunctor{\mathrm{Quot}}
\def\Hilbfunctor{\mathrm{Hilb}}
\def\Curvesstack{\mathcal{C}\!\mathit{urves}}
\def\Polarizedstack{\mathcal{P}\!\mathit{olarized}}
\def\Complexesstack{\mathcal{C}\!\mathit{omplexes}}
\def\Pic{\mathop{\mathrm{Pic}}\nolimits}
\def\Picardstack{\mathcal{P}\!\mathit{ic}}
\def\Picardfunctor{\mathrm{Pic}}
\def\Deformationcategory{\mathcal{D}\!\mathit{ef}}
\usepackage{cleveref}
\crefname{originaltheorem}{original theorem}{original theorems}
\Crefname{originaltheorem}{Original theorem}{Original theorems}
\makeatletter
\newlabel{groupoids-lemma-etale-equivalence-relation}{{0DT7}{source}{}{sourcepointer-1}{}}
\newlabel{groupoids-lemma-etale-equivalence-relation@cref}{{[originaltheorem][0][]0DT7}{[1][0][]source}{}{}{}}
\makeatother
\crefname{originaltheorem}{original theorem}{original theorems}
\Crefname{originaltheorem}{Original theorem}{Original theorems}
\makeatletter
\newlabel{lemma-etale-local-finite-inertia}{{lemma-etale-local-finite-inertia}{source}{}{sourcepointer-2}{}}
\newlabel{lemma-etale-local-finite-inertia@cref}{{[originaltheorem][0][]lemma-etale-local-finite-inertia}{[1][0][]source}{}{}{}}
\makeatother
\crefname{originaltheorem}{original theorem}{original theorems}
\Crefname{originaltheorem}{Original theorem}{Original theorems}
\makeatletter
\newlabel{lemma-etale-separated-over-well-nigh-affine}{{lemma-etale-separated-over-well-nigh-affine}{source}{}{sourcepointer-3}{}}
\newlabel{lemma-etale-separated-over-well-nigh-affine@cref}{{[originaltheorem][0][]lemma-etale-separated-over-well-nigh-affine}{[1][0][]source}{}{}{}}
\makeatother
\crefname{originaltheorem}{original theorem}{original theorems}
\Crefname{originaltheorem}{Original theorem}{Original theorems}
\makeatletter
\newlabel{lemma-moduli-space-finite-affine}{{lemma-moduli-space-finite-affine}{source}{}{sourcepointer-4}{}}
\newlabel{lemma-moduli-space-finite-affine@cref}{{[originaltheorem][0][]lemma-moduli-space-finite-affine}{[1][0][]source}{}{}{}}
\makeatother
\crefname{originaltheorem}{original theorem}{original theorems}
\Crefname{originaltheorem}{Original theorem}{Original theorems}
\makeatletter
\newlabel{lemma-well-nigh-affine-moduli-space}{{lemma-well-nigh-affine-moduli-space}{source}{}{sourcepointer-5}{}}
\newlabel{lemma-well-nigh-affine-moduli-space@cref}{{[originaltheorem][0][]lemma-well-nigh-affine-moduli-space}{[1][0][]source}{}{}{}}
\makeatother
\crefname{originaltheorem}{original theorem}{original theorems}
\Crefname{originaltheorem}{Original theorem}{Original theorems}
\makeatletter
\newlabel{spaces-theorem-presentation}{{02WW}{source}{}{sourcepointer-6}{}}
\newlabel{spaces-theorem-presentation@cref}{{[originaltheorem][0][]02WW}{[1][0][]source}{}{}{}}
\makeatother
\crefname{originaltheorem}{original theorem}{original theorems}
\Crefname{originaltheorem}{Original theorem}{Original theorems}
\makeatletter
\newlabel{stacks-morphisms-lemma-base-change-etale}{{0CIN}{source}{}{sourcepointer-7}{}}
\newlabel{stacks-morphisms-lemma-base-change-etale@cref}{{[originaltheorem][0][]0CIN}{[1][0][]source}{}{}{}}
\makeatother
\crefname{originaltheorem}{original theorem}{original theorems}
\Crefname{originaltheorem}{Original theorem}{Original theorems}
\makeatletter
\newlabel{stacks-morphisms-lemma-base-change-separated}{{050F}{source}{}{sourcepointer-8}{}}
\newlabel{stacks-morphisms-lemma-base-change-separated@cref}{{[originaltheorem][0][]050F}{[1][0][]source}{}{}{}}
\makeatother
\crefname{originaltheorem}{original theorem}{original theorems}
\Crefname{originaltheorem}{Original theorem}{Original theorems}
\makeatletter
\newlabel{stacks-morphisms-lemma-base-change-stabilizer-preserving}{{0DUB}{source}{}{sourcepointer-9}{}}
\newlabel{stacks-morphisms-lemma-base-change-stabilizer-preserving@cref}{{[originaltheorem][0][]0DUB}{[1][0][]source}{}{}{}}
\makeatother
\crefname{originaltheorem}{original theorem}{original theorems}
\Crefname{originaltheorem}{Original theorem}{Original theorems}
\makeatletter
\newlabel{stacks-morphisms-lemma-check-quasi-compact-covering}{{0DTL}{source}{}{sourcepointer-10}{}}
\newlabel{stacks-morphisms-lemma-check-quasi-compact-covering@cref}{{[originaltheorem][0][]0DTL}{[1][0][]source}{}{}{}}
\makeatother
\crefname{originaltheorem}{original theorem}{original theorems}
\Crefname{originaltheorem}{Original theorem}{Original theorems}
\makeatletter
\newlabel{stacks-morphisms-lemma-check-separated-covering}{{06TZ}{source}{}{sourcepointer-11}{}}
\newlabel{stacks-morphisms-lemma-check-separated-covering@cref}{{[originaltheorem][0][]06TZ}{[1][0][]source}{}{}{}}
\makeatother
\crefname{originaltheorem}{original theorem}{original theorems}
\Crefname{originaltheorem}{Original theorem}{Original theorems}
\makeatletter
\newlabel{stacks-morphisms-lemma-check-universal-homeomorphism-covering}{{0DTQ}{source}{}{sourcepointer-12}{}}
\newlabel{stacks-morphisms-lemma-check-universal-homeomorphism-covering@cref}{{[originaltheorem][0][]0DTQ}{[1][0][]source}{}{}{}}
\makeatother
\crefname{originaltheorem}{original theorem}{original theorems}
\Crefname{originaltheorem}{Original theorem}{Original theorems}
\makeatletter
\newlabel{stacks-morphisms-lemma-compose-after-separated}{{050M}{source}{}{sourcepointer-13}{}}
\newlabel{stacks-morphisms-lemma-compose-after-separated@cref}{{[originaltheorem][0][]050M}{[1][0][]source}{}{}{}}
\makeatother
\crefname{originaltheorem}{original theorem}{original theorems}
\Crefname{originaltheorem}{Original theorem}{Original theorems}
\makeatletter
\newlabel{stacks-morphisms-lemma-etale-iso}{{0DUD}{source}{}{sourcepointer-14}{}}
\newlabel{stacks-morphisms-lemma-etale-iso@cref}{{[originaltheorem][0][]0DUD}{[1][0][]source}{}{}{}}
\makeatother
\crefname{originaltheorem}{original theorem}{original theorems}
\Crefname{originaltheorem}{Original theorem}{Original theorems}
\makeatletter
\newlabel{stacks-morphisms-lemma-etale-permanence}{{0CIR}{source}{}{sourcepointer-15}{}}
\newlabel{stacks-morphisms-lemma-etale-permanence@cref}{{[originaltheorem][0][]0CIR}{[1][0][]source}{}{}{}}
\makeatother
\crefname{originaltheorem}{original theorem}{original theorems}
\Crefname{originaltheorem}{Original theorem}{Original theorems}
\makeatletter
\newlabel{stacks-morphisms-lemma-universally-injective}{{0CPN}{source}{}{sourcepointer-16}{}}
\newlabel{stacks-morphisms-lemma-universally-injective@cref}{{[originaltheorem][0][]0CPN}{[1][0][]source}{}{}{}}
\makeatother
\crefname{originaltheorem}{original theorem}{original theorems}
\Crefname{originaltheorem}{Original theorem}{Original theorems}
\makeatletter
\newlabel{stacks-morphisms-lemma-universally-injective-point}{{0DTM}{source}{}{sourcepointer-17}{}}
\newlabel{stacks-morphisms-lemma-universally-injective-point@cref}{{[originaltheorem][0][]0DTM}{[1][0][]source}{}{}{}}
\makeatother
\crefname{originaltheorem}{original theorem}{original theorems}
\Crefname{originaltheorem}{Original theorem}{Original theorems}
\makeatletter
\newlabel{stacks-morphisms-remark-identify-automorphism-groups}{{0DTW}{source}{}{sourcepointer-18}{}}
\newlabel{stacks-morphisms-remark-identify-automorphism-groups@cref}{{[originaltheorem][0][]0DTW}{[1][0][]source}{}{}{}}
\makeatother
\title{Keel-Mori coarse moduli space for finite-inertia stacks}
\author{The Stacks Project Authors}
\date{stacks-project:stacks-more-morphisms.tex:local-snapshot-20261012; GNU FDL 1.2 or later}
\begin{document}
\maketitle
\noindent\textit{Editorial note.} Statement span is lines 3853-3862 of sources/stacks-stacks-more-morphisms/pinned-original-source.tex; proof span is lines 3864-4102, the complete author proof (237 body lines).\par\medskip
\noindent\textit{Editorial note.} Navigation stubs: 18 cited-but-not-included labels are supplied as non-load-bearing pointers with their pinned tags where the snapshot tag index has them.\par\medskip
\noindent\textit{Editorial note.} extra\_packages is empty on purpose: the retained preamble already loads xy with options and a second declaration clashes in the assembler wrapper.\par\medskip
\noindent\textit{Editorial note.} Text is retained verbatim, including the author's own typos and placeholders.\par\medskip
\noindent\textit{Editorial note.} Queue index 0 of queue\_stacks-stacks-more-morphisms.json (pb=9173, stmt 3853-3862, proof 3864-4102).\par\medskip
\noindent\textit{Editorial note.} stacks-stacks-more-morphisms wave smm01 (2026-10-14): fresh adjudication accept per lane-s3/smm\_judgments.md; pool per Chief 10:5x open-pool ruling (net 59 with-proof rows, bloodline lane-s3 by anchor audit, v2 true-key namespace stacks:stacks-more-morphisms.tex, row lemma-deform-property-fp-over-ft avoided on cross-chapter key collision); 60 with-proof rows judged 45 accept / 14 reject / 1 avoid; patch-sd-envs contract spans scanned.\par\medskip
\noindent\textit{Editorial note.} Cross-chapter same-label warning (NOT an occupancy conflict, Chief ruling 2026-10-13 evening): label 'theorem-keel-mori' also exists in stacks-snapshot-20261011 — a different proposition in a different chapter file.\par\medskip
\noindent\textit{Allowed standard prerequisites (editorial).}\begin{itemize}
\item Explicit prerequisite boundary. The cited inputs are prerequisites; the argument in the retained span is not replaced by any cited result.
\end{itemize}\medskip
\ifdefined\hypertarget\hypertarget{sourcepointer-1}{}\fi
\noindent\textit{Source reference.} Original theorem 0DT7: sources/stacks-stacks-more-morphisms/pinned-original-source.tex, source lines 4102--4102. Cited by the retained statement or proof but not part of the collected spans; supplied only as a navigation pointer. The argument does not depend on its content.\par\medskip
\ifdefined\hypertarget\hypertarget{sourcepointer-2}{}\fi
\noindent\textit{Source reference.} Original theorem lemma-etale-local-finite-inertia: sources/stacks-stacks-more-morphisms/pinned-original-source.tex, source lines 3798--3814. Cited by the retained statement or proof but not part of the collected spans; supplied only as a navigation pointer. The argument does not depend on its content.\par\medskip
\ifdefined\hypertarget\hypertarget{sourcepointer-3}{}\fi
\noindent\textit{Source reference.} Original theorem lemma-etale-separated-over-well-nigh-affine: sources/stacks-stacks-more-morphisms/pinned-original-source.tex, source lines 3692--3708. Cited by the retained statement or proof but not part of the collected spans; supplied only as a navigation pointer. The argument does not depend on its content.\par\medskip
\ifdefined\hypertarget\hypertarget{sourcepointer-4}{}\fi
\noindent\textit{Source reference.} Original theorem lemma-moduli-space-finite-affine: sources/stacks-stacks-more-morphisms/pinned-original-source.tex, source lines 3592--3599. Cited by the retained statement or proof but not part of the collected spans; supplied only as a navigation pointer. The argument does not depend on its content.\par\medskip
\ifdefined\hypertarget\hypertarget{sourcepointer-5}{}\fi
\noindent\textit{Source reference.} Original theorem lemma-well-nigh-affine-moduli-space: sources/stacks-stacks-more-morphisms/pinned-original-source.tex, source lines 3425--3433. Cited by the retained statement or proof but not part of the collected spans; supplied only as a navigation pointer. The argument does not depend on its content.\par\medskip
\ifdefined\hypertarget\hypertarget{sourcepointer-6}{}\fi
\noindent\textit{Source reference.} Original theorem 02WW: sources/stacks-stacks-more-morphisms/pinned-original-source.tex, source lines 4102--4102. Cited by the retained statement or proof but not part of the collected spans; supplied only as a navigation pointer. The argument does not depend on its content.\par\medskip
\ifdefined\hypertarget\hypertarget{sourcepointer-7}{}\fi
\noindent\textit{Source reference.} Original theorem 0CIN: sources/stacks-stacks-more-morphisms/pinned-original-source.tex, source lines 4102--4102. Cited by the retained statement or proof but not part of the collected spans; supplied only as a navigation pointer. The argument does not depend on its content.\par\medskip
\ifdefined\hypertarget\hypertarget{sourcepointer-8}{}\fi
\noindent\textit{Source reference.} Original theorem 050F: sources/stacks-stacks-more-morphisms/pinned-original-source.tex, source lines 4102--4102. Cited by the retained statement or proof but not part of the collected spans; supplied only as a navigation pointer. The argument does not depend on its content.\par\medskip
\ifdefined\hypertarget\hypertarget{sourcepointer-9}{}\fi
\noindent\textit{Source reference.} Original theorem 0DUB: sources/stacks-stacks-more-morphisms/pinned-original-source.tex, source lines 4102--4102. Cited by the retained statement or proof but not part of the collected spans; supplied only as a navigation pointer. The argument does not depend on its content.\par\medskip
\ifdefined\hypertarget\hypertarget{sourcepointer-10}{}\fi
\noindent\textit{Source reference.} Original theorem 0DTL: sources/stacks-stacks-more-morphisms/pinned-original-source.tex, source lines 4102--4102. Cited by the retained statement or proof but not part of the collected spans; supplied only as a navigation pointer. The argument does not depend on its content.\par\medskip
\ifdefined\hypertarget\hypertarget{sourcepointer-11}{}\fi
\noindent\textit{Source reference.} Original theorem 06TZ: sources/stacks-stacks-more-morphisms/pinned-original-source.tex, source lines 4102--4102. Cited by the retained statement or proof but not part of the collected spans; supplied only as a navigation pointer. The argument does not depend on its content.\par\medskip
\ifdefined\hypertarget\hypertarget{sourcepointer-12}{}\fi
\noindent\textit{Source reference.} Original theorem 0DTQ: sources/stacks-stacks-more-morphisms/pinned-original-source.tex, source lines 4102--4102. Cited by the retained statement or proof but not part of the collected spans; supplied only as a navigation pointer. The argument does not depend on its content.\par\medskip
\ifdefined\hypertarget\hypertarget{sourcepointer-13}{}\fi
\noindent\textit{Source reference.} Original theorem 050M: sources/stacks-stacks-more-morphisms/pinned-original-source.tex, source lines 4102--4102. Cited by the retained statement or proof but not part of the collected spans; supplied only as a navigation pointer. The argument does not depend on its content.\par\medskip
\ifdefined\hypertarget\hypertarget{sourcepointer-14}{}\fi
\noindent\textit{Source reference.} Original theorem 0DUD: sources/stacks-stacks-more-morphisms/pinned-original-source.tex, source lines 4102--4102. Cited by the retained statement or proof but not part of the collected spans; supplied only as a navigation pointer. The argument does not depend on its content.\par\medskip
\ifdefined\hypertarget\hypertarget{sourcepointer-15}{}\fi
\noindent\textit{Source reference.} Original theorem 0CIR: sources/stacks-stacks-more-morphisms/pinned-original-source.tex, source lines 4102--4102. Cited by the retained statement or proof but not part of the collected spans; supplied only as a navigation pointer. The argument does not depend on its content.\par\medskip
\ifdefined\hypertarget\hypertarget{sourcepointer-16}{}\fi
\noindent\textit{Source reference.} Original theorem 0CPN: sources/stacks-stacks-more-morphisms/pinned-original-source.tex, source lines 4102--4102. Cited by the retained statement or proof but not part of the collected spans; supplied only as a navigation pointer. The argument does not depend on its content.\par\medskip
\ifdefined\hypertarget\hypertarget{sourcepointer-17}{}\fi
\noindent\textit{Source reference.} Original theorem 0DTM: sources/stacks-stacks-more-morphisms/pinned-original-source.tex, source lines 4102--4102. Cited by the retained statement or proof but not part of the collected spans; supplied only as a navigation pointer. The argument does not depend on its content.\par\medskip
\ifdefined\hypertarget\hypertarget{sourcepointer-18}{}\fi
\noindent\textit{Source reference.} Original theorem 0DTW: sources/stacks-stacks-more-morphisms/pinned-original-source.tex, source lines 4102--4102. Cited by the retained statement or proof but not part of the collected spans; supplied only as a navigation pointer. The argument does not depend on its content.\par\medskip
\begin{theorem}[Keel-Mori]
\label{theorem-keel-mori}
Let $\mathcal{X}$ be an algebraic stack. Assume
$\mathcal{I}_\mathcal{X} \to \mathcal{X}$ is finite.
Then there exists a uniform categorical moduli space
$$
f : \mathcal{X} \longrightarrow M
$$
and $f$ is separated, quasi-compact, and a universal homeomorphism.
\end{theorem}

\begin{proof}
We choose a set $I$\footnote{The reader who is still keeping
track of set theoretic issues should make sure $I$ is not too large.}
and for $i \in I$ a morphism of algebraic stacks
$g_i : \mathcal{X}_i \to \mathcal{X}$ as in
Lemma \ref{lemma-etale-local-finite-inertia}; we will use all
of the properties listed in this lemma without further mention.
Let
$$
f_i : \mathcal{X}_i \to M_i
$$
be as in Lemma \ref{lemma-well-nigh-affine-moduli-space}.
Consider the stacks
$$
\mathcal{X}_{ij} = \mathcal{X}_i \times_{g_i, \mathcal{X}, g_j} \mathcal{X}_j
$$
for $i, j \in I$. The projections $\mathcal{X}_{ij} \to \mathcal{X}_i$
and $\mathcal{X}_{ij} \to \mathcal{X}_j$ are separated
by Morphisms of Stacks, Lemma
\ref{stacks-morphisms-lemma-base-change-separated},
\'etale by Morphisms of Stacks, Lemma
\ref{stacks-morphisms-lemma-base-change-etale},
and induce isomorphisms on automorphism groups
(as in Morphisms of Stacks, Remark
\ref{stacks-morphisms-remark-identify-automorphism-groups}) by
Morphisms of Stacks, Lemma
\ref{stacks-morphisms-lemma-base-change-stabilizer-preserving}.
Thus we may apply Lemma \ref{lemma-etale-separated-over-well-nigh-affine}
to find a commutative diagram
$$
\xymatrix{
\mathcal{X}_i \ar[d]_{f_i} &
\mathcal{X}_{ij} \ar[d]_{f_{ij}} \ar[l] \ar[r] &
\mathcal{X}_j \ar[d]_{f_j} \\
M_i &
M_{ij} \ar[l] \ar[r] &
M_j
}
$$
with cartesian squares where $M_{ij} \to M_i$ and $M_{ij} \to M_j$
are separated \'etale morphisms of schemes; here we also use that $f_i$
is a uniform categorical quotient by
Lemma \ref{lemma-moduli-space-finite-affine}.
Claim:
$$
\coprod M_{ij} \longrightarrow \coprod M_i \times \coprod M_i
$$
is an \'etale equivalence relation.

\medskip\noindent
Proof of the claim. Set $R = \coprod M_{ij}$ and $U = \coprod M_i$.
We have already seen that $t : R \to U$ and $s : R \to U$ are \'etale.
Let us construct a morphism $c : R \times_{s, U, t} R \to R$
compatible with $\text{pr}_{13} : U \times U \times U \to U \times U$.
Namely, for $i, j, k \in I$ we consider
$$
\mathcal{X}_{ijk} =
\mathcal{X}_i \times_{g_i, \mathcal{X}, g_j} \mathcal{X}_j
\times_{g_j, \mathcal{X}, g_k} \mathcal{X}_k =
\mathcal{X}_{ij} \times_{\mathcal{X}_j} \mathcal{X}_{jk}
$$
Arguing exactly as in the previous paragraph,
we find that $M_{ijk} = M_{ij} \times_{M_j} M_{jk}$
is a categorical moduli space for $\mathcal{X}_{ijk}$.
In particular, there is a canonical morphism
$M_{ijk} = M_{ij} \times_{M_j} M_{jk} \to M_{ik}$
coming from the projection $\mathcal{X}_{ijk} \to \mathcal{X}_{ik}$.
Putting these morphisms together we obtain the morphism $c$.
In a similar fashion we construct a morphism $e : U \to R$
compatible with $\Delta : U \to U \times U$ and
$i : R \to R$ compatible with the flip $U \times U \to U \times U$.
Let $k$ be an algebraically closed field. Then
$$
\Mor(\Spec(k), \mathcal{X}_i) \to \Mor(\Spec(k), M_i) = M_i(k)
$$
is bijective on isomorphism classes and the same remains true after any
base change by a morphism $M' \to M$. This follows from our choice
of $f_i$ and Morphisms of Stacks, Lemmas
\ref{stacks-morphisms-lemma-universally-injective} and
\ref{stacks-morphisms-lemma-universally-injective-point}.
By construction of $2$-fibred products the diagram
$$
\xymatrix{
\Mor(\Spec(k), \mathcal{X}_{ij}) \ar[d] \ar[r] &
\Mor(\Spec(k), \mathcal{X}_j) \ar[d] \\
\Mor(\Spec(k), \mathcal{X}_i) \ar[r] &
\Mor(\Spec(k), \mathcal{X})
}
$$
is a fibre product of categories. By our choice of $g_i$ the
functors in this diagram induce bijections on automorphism groups.
It follows that this diagram induces a fibre product diagram
on sets of isomorphism classes! Thus we see that
$$
R(k) = U(k) \times_{|\Mor(\Spec(k), \mathcal{X})|} U(k)
$$
where $|\Mor(\Spec(k), \mathcal{X})|$ denotes the set
of isomorphism classes.
In particular, for any algebraically closed field $k$
the map on $k$-valued point is an equivalence relation.
We conclude the claim holds by
Groupoids, Lemma \ref{groupoids-lemma-etale-equivalence-relation}.

\medskip\noindent
Let $M = U/R$ be the algebraic space which is the quotient of the above
\'etale equivalence relation, see
Spaces, Theorem \ref{spaces-theorem-presentation}.
There is a canonical morphism $f : \mathcal{X} \to M$
fitting into commutative diagrams
\begin{equation}
\label{equation-fundamental-diagram}
\xymatrix{
\mathcal{X}_i \ar[r]_{g_i} \ar[d]_{f_i} & \mathcal{X} \ar[d]^f \\
M_i \ar[r] & M
}
\end{equation}
Namely, such a morphism $f$ is given by a functor
$$
f : \Mor(T, \mathcal{X}) \longrightarrow \Mor(T, M)
$$
for any scheme $T$ compatible with base change. Let $a : T \to \mathcal{X}$
be an object of the left hand side. We obtain an \'etale covering
$\{T_i \to T\}$ with $T_i = \mathcal{X}_i \times_\mathcal{X} T$
and morphisms $a_i : T_i \to \mathcal{X}_i$. Then we get
$b_i = f_i \circ a_i : T_i \to M_i$. Since
$T_i \times_T T_j = \mathcal{X}_{ij} \times_\mathcal{X} T$
we moreover get a morphism $a_{ij} : T_i \times_T T_j \to \mathcal{X}_{ij}$.
Setting $b_{ij} = f_{ij} \circ a_{ij}$ we find that
$b_i \times b_j$ factors through the monomorphism
$M_{ij} \to M_i \times M_j$. Hence the morphisms
$$
T_i \xrightarrow{b_i} M_i \to M
$$
agree on $T_i \times_T T_j$. As $M$ is a sheaf for the \'etale
topology, we see that these morphisms glue to a unique
morphism $b = f(a) : T \to M$. We omit the verification that
this construction is compatible with base change and we omit
the verification that the diagrams
(\ref{equation-fundamental-diagram}) commute.

\medskip\noindent
Claim: the diagrams (\ref{equation-fundamental-diagram}) are cartesian.
To see this we study the induced morphism
$$
h_i : \mathcal{X}_i \longrightarrow M_i \times_M \mathcal{X}
$$
This is a morphism of stacks \'etale over $\mathcal{X}$
and hence $h_i$ is \'etale (Morphisms of Stacks, Lemma
\ref{stacks-morphisms-lemma-etale-permanence}).
Since $g_i$ is separated, we see $h_i$ is separated
(use Morphisms of Stacks, Lemma
\ref{stacks-morphisms-lemma-compose-after-separated} and the fact
seen above that the diagonal of $\mathcal{X}$ is separated).
The morphism $h_i$ induces isomorphisms on automorphism groups
(Morphisms of Stacks, Remark
\ref{stacks-morphisms-remark-identify-automorphism-groups})
as this is true for $g_i$. For an algebraically closed field $k$
the diagram
$$
\xymatrix{
\Mor(\Spec(k), M_i \times_M \mathcal{X}) \ar[r] \ar[d] &
\Mor(\Spec(k), \mathcal{X}) \ar[d] \\
M_i(k) \ar[r] &
M(k)
}
$$
is a catesian diagram of categories and the top arrow
induces bijections on automorphism groups.
On the other hand, we have
$$
M(k) = U(k)/R(k) = U(k)/
U(k) \times_{|\Mor(\Spec(k), \mathcal{X})|} U(k) =
|\Mor(\Spec(k), \mathcal{X})|
$$
by what we said above. Thus the right vertical arrow in the
cartesian diagram above is a bijection on isomorphism classes.
We conclude that
$|\Mor(\Spec(k), M_i \times_M \mathcal{X})| \to M_i(k)$ is bijective.
Review: $h_i$ is a separated, \'etale, induces isomorphisms on
automorphism groups (as in Morphisms of Stacks, Remark
\ref{stacks-morphisms-remark-identify-automorphism-groups}), and
induces an equivalence on fibre categories over algebraically closed fields.
Hence it is an isomorphism by Morphisms of Stacks, Lemma
\ref{stacks-morphisms-lemma-etale-iso}.

\medskip\noindent
From the claim we get in particular the following:
we have a surjective \'etale morphism $U \to M$
such that the base change of $f$ is separated, quasi-compact,
and a universal homeomorphism. It follows that $f$ is separated,
quasi-compact, and a universal homeomorphism.
See Morphisms of Stacks, Lemma
\ref{stacks-morphisms-lemma-check-separated-covering},
\ref{stacks-morphisms-lemma-check-quasi-compact-covering}, and
\ref{stacks-morphisms-lemma-check-universal-homeomorphism-covering}

\medskip\noindent
To finish the proof we have to show that $f : \mathcal{X} \to M$
is a uniform categorical moduli space.
To prove this it suffices to show that given a flat morphism
$M' \to M$ of algebraic spaces, the base change
$$
M' \times_M \mathcal{X} \longrightarrow M'
$$
is a categorical moduli space. Thus we consider a morphism
$$
\theta : M' \times_M \mathcal{X} \longrightarrow E
$$
where $E$ is an algebraic space. For each $i$ we know that
$f_i$ is a uniform categorical moduli space. Hence we obtain
$$
\xymatrix{
M' \times_M \mathcal{X}_i \ar[d] \ar[r] &
M' \times_M \mathcal{X} \ar[d]^\theta \\
M' \times_M M_i \ar[r]^{\psi_i} &
E
}
$$
Since $\{M' \times_M M_i \to M'\}$ is an \'etale covering,
to obtain the desired morphism $\psi : M' \to E$ it suffices
to show that $\psi_i$ and $\psi_j$ agree over
$M' \times_M M_i \times_M M_j = M' \times_M M_{ij}$.
This follows easily from the fact that
$f_{ij} : \mathcal{X}_{ij} =
\mathcal{X}_i \times_\mathcal{X} \mathcal{X}_j \to M_{ij}$ is a uniform
categorical quotient; details omitted.
Then finally one shows that $\psi$ fits into the commutative diagram
$$
\xymatrix{
M' \times_M \mathcal{X} \ar[d] \ar[rd]^\theta \\
M' \ar[r]^\psi &
E
}
$$
because ``$\{M' \times_M \mathcal{X}_i \to M' \times_M \mathcal{X}\}$
is an \'etale covering'' and the morphisms $\psi_i$ fit into the
corresponding commutative diagrams by construction.
This finishes the proof of the Keel-Mori theorem.
\end{proof}

\end{document}
